% THRESHOLD-DESCENT REPORT::THRES-HASHBRIDGE4::2026-10-18_1416_KST
\documentclass{article}
\usepackage{amsmath,amssymb,amsthm}
\newtheorem{theorem}{Theorem}
\title{THRES-HASHBRIDGE4: Lossy Hash Fingerprints and the MCSP Promise}
\date{2026-10-18}\begin{document}\maketitle
\noindent STATUS: PROVED/EXACT restricted lemmas pending independent review; COMPUTATIONAL tests; threshold descent and P versus NP OPEN.

\section{IDEAS PROPOSED BEFORE LITERATURE}
\paragraph{A.} Test dense linear sketches via hard vectors in the kernel.
\paragraph{B.} Use polynomial fingerprints as a randomized lossy semantic MCSP quotient.
\paragraph{C.} Test whether hash bits themselves can be treated as a smaller-MCSP truth table.
\paragraph{D.} Audit direct sums, restrictions, padding, local PRGs, and implicit encodings for source promise, target model and slope preservation.

\section{LITERATURE CHECK}
Primary sources: CHMY STACS 2021 exact one-tape MCSP lower bound, MMW STOC 2019 search streaming upper bound, OPS Theory of Computing 2021 Gap-MCSP circuit magnification, Carter--Wegman 1979 universal hashing, Karp--Rabin 1987 polynomial fingerprints. Reviewed project ASYMPTOTIC-CONE5, MCSPMAGCONE1, THRES-ENTROPY1, THRES-LOSSYPROBE2, THRES-SHANNON3 local source, MAGRED-QCOVER2 and MAGRED-TAPEFUSION3. Earlier entropy, input-query, gap-slice and one-tape fusion restrictions are prior, not newly established here.

\section{CONSTRUCTION AND PROOF}
Write $N=2^n$, $s=N^\beta$ with fixed $0<\beta<1$, and let $Y_s$ be the set of $N$-bit truth tables with fan-in-two circuit size at most $s$. A syntactic count gives $|Y_s|\le U=(s+1)(n+s+2)[3(n+s+2)^2]^s$; put $B=\lceil\log_2 U\rceil=o(N)$.

\begin{theorem}[AFFINE-KERNEL4]
Any fixed affine sketch $f\mapsto Af+u$ followed by an arbitrary deterministic decoder fails to decide MCSP[$s$] exactly if $\dim\ker A>B$. In particular an output of fewer than $N-B$ bits cannot work.
\end{theorem}
\noindent Proof: The kernel has more than $|Y_s|$ members, hence contains a hard NO word as well as zero YES, all mapping to $u$.

Choose $b=B+\lceil\log_2N\rceil+3$, a field $\mathbb F_{2^b}$, and $H_\alpha(f)=\sum_{i=0}^{N-1}f_i\alpha^i$. Define the indexed image language
$\mathcal I_{n,s,b}=\{(P,\alpha,z):\exists \text{ size-}s\text{ circuit } C,\ H_\alpha(TT(C))=z\}$, where $P$ encodes the field.

\begin{theorem}[FINGERPRINT-QUOTIENT4]
For each fixed source $f$, with $\alpha$ random uniform, the target $(P,\alpha,H_\alpha(f))$ has length $N^{\beta+o(1)}$ and belongs to $\mathcal I$ with probability one if $C(f)\le s$, and with probability at most $1/8$ if $C(f)>s$.
\end{theorem}
\noindent Proof: YES has its own circuit witness. For NO and each YES $g$, the nonzero polynomial $\sum_i(f_i-g_i)z^i$ has at most $N-1$ roots. The union bound is $U(N-1)/2^b\le1/8$. Quantifiers are $(\forall f)\Pr_\alpha[\text{correct}]\ge7/8$, not $(\exists\alpha)(\forall f)$.

\begin{theorem}[FINGERPRINT-YES-FLIP4]
Let $b=2^m$, $r=b^\gamma$, $0<\gamma<1$. The easy source point function $f=e_1$ has $H_\alpha(f)=\alpha$, a uniformly random $b$-bit string. When interpreted as an ordinary target MCSP[$r$] truth table,
$\Pr_\alpha[C(H_\alpha(f))\le r]\le2^{-b+O(r\log(m+r))}=2^{-b+o(b)}$.
\end{theorem}
\noindent Proof: a random $b$-bit table has probability at most $U(m,r)/2^b$ of being generated by a circuit with $r$ gates. Random Toeplitz sketches have the same YES-to-NO behavior on any easy nonzero point table: the relevant column consists of independent random diagonal bits.

\section{TESTS AND OUTCOME}
The full local audit performed 32,640 GF($2^4$) polynomial-pair checks across 522,240 evaluations, all meeting the degree-root bound. Every nonzero four-bit input under all 2,048 Toeplitz seeds yielded exactly uniform eight-bit outputs, checking 3,840 histogram bins. Exhaustive three-variable circuit enumeration gave 5, 14 and 40 easy tables of 256 at gate budgets 0, 1 and 2.

The quotient target is an indexed hash-image language, NOT ordinary MCSP or OPS Gap-MCSP. Its explicit witness verifier scans $N$ source positions: for each fixed $\beta$ polynomial target length, but as $\beta\downarrow0$ the displayed verifier's exponent $1/\beta$ is not fixed. The straightforward bit-serial fingerprint requires at least $N b=N^{1+\beta+o(1)}$ elementary updates, exceeding CHMY high-parameter lower exponent $1+g$ with $g<2\beta-1<\beta$; no general hashing lower bound is asserted. Therefore no valid $(\lambda,\rho)$, slope increase, or P versus NP proof follows.

NEXT: HASH-IMAGE-BRIDGE1, a uniformly checkable, fully charged, same-one-tape compiler from image membership to the exact ordinary MCSP or specified Gap-MCSP promise, preserving a common exponent for $\beta\to0$. Root main.tex unchanged due absent specialist ownership.
\end{document}
% END_OF_VERIFIED_REPORT::THRESHOLD-DESCENT::THRES-HASHBRIDGE4::P_VS_NP_OPEN